\documentclass[11pt]{article}
\usepackage[a4paper,margin=25mm]{geometry}
\usepackage{amsmath,amssymb,amsthm}
\usepackage[hidelinks]{hyperref}
\newtheorem{theorem}{Theorem}
\begin{document}
\title{A counterexample to the unrestricted bound in Conjecture 3464}
\author{Chengzhe Gao}\date{}\maketitle
The English and Chinese versions both assert the universal bound
$\lambda(G)\leq\Delta(G)^2$, without a lower bound on $\Delta$.
Here $\lambda$ is the least label range for an $L(2,1)$-labeling.
This note disproves that unrestricted assertion. It does not disprove
the classical version restricted to $\Delta\geq2$.

\begin{theorem}
For the single-edge graph $K_2$, one has $\Delta(K_2)=1$ and
$\lambda(K_2)=2>1=\Delta(K_2)^2$.
\end{theorem}
\begin{proof}
Each vertex has exactly one neighbor. In any valid labeling, the two
adjacent vertices have labels differing by at least two. Hence their
label range is at least two, regardless of how large the labels are.
The labels zero and two attain range two. There are no vertices at
distance exactly two, so the second distance condition is vacuous.
\end{proof}

The later assertions in the source form part of a conjunction containing
the displayed universal bound; a counterexample to this necessary part
already disproves the conjecture as stated. No asymptotic claim is used.
The omitted degree restriction is substantial: the standard formulation
explicitly assumes $\Delta\geq2$; see the statement in
\href{https://doi.org/10.1016/j.dam.2025.10.011}{Discrete Applied Mathematics,
DOI 10.1016/j.dam.2025.10.011}.

\paragraph{Formal verification.}
The Lean file defines finite simple graphs with symmetry and no loops,
computes degrees from the entire vertex list, and defines distance two
as nonadjacency together with a common neighbor. The definition
\texttt{L21} imposes both standard separation conditions.
\texttt{SpanAtMost f R} means that every label lies in some translated
interval $[\ell,\ell+R]$. The theorem \texttt{span\_range\_equivalence}
proves that this agrees with maximum minus minimum on this graph.
\texttt{every\_valid\_span\_at\_least\_two} quantifies over all natural
labels and all intervals; it is not a search with an imposed label bound.
The labels zero and two and this universal lower bound yield
\texttt{edge\_minimum\_span}. Finally,
\texttt{conjecture3464\_false} negates the claimed universal bound itself.
Integer labels give the same minimum range by translating a finite
labeling so that its minimum is zero.

The program uses Lean 4.19.0 and \texttt{Std}. The printed axiom lists
contain only Lean's standard logical foundations; no unproved declaration,
external oracle, or native evaluation tactic occurs.

\end{document}
